\section{Pokrycie kart i mechanik gry}

Jednym z ważnych parametrów symulatora gry karcianej jest pokrycie puli dostępnych kart. Niepełny zakres implementacji kart ogranicza przestrzeń badawczą oraz może wpływać na wiarygodność wyników eksperymentów. SabberStone pokrywa ponad 98\% kart wchodzących w skład formatu standardowego aktywnego w 2019 roku, osiągając 100\% pokrycie dla wszystkich podstawowych i klasycznych zestawów kart oraz starszych dodatków. Karty z nowszych rozszerzeń, takich jak \textit{Rise of Shadows} czy \textit{The Boomsday Project}, zaimplementowane są w przedziale 94--99\% \cite{sabberstone_github}.

Mechanizm implementacji efektów kart jest oparty na deklaratywnym systemie zadań, w którym każdy efekt, taki jak zadawanie obrażeń, dobieranie kart czy przywoływanie jednostek, stanowi odrębny komponent. Złożone efekty konstruuje się poprzez dodawanie tych elementów do siebie, co znacząco ułatwia implementację kart z kolejnych rozszerzeń. Osobnym wyzwaniem jest obsługa dynamicznych modyfikatorów statystyk wynikających z wzmocnień i aur. Silnik radzi sobie z tymi elementami za pomocą autorskiego systemu warstwowego, który oblicza finalne wartości atrybutów jednostek jako złożenie wszystkich aktywnych efektów, zgodnie z oficjalną hierarchią reguł gry \cite{sabberstone_hearthsim}.
